\subsection{The marked-line intersection and a finite higher obstruction}

\begin{proposition}[D-028: \status{Proved}; an exact reduction and a higher image class]
\label{sa08:D028}
Let \(D_N\) be the differential image of
Statement~\ref{sa08:D027}, over a commutative ring \(k\).
For a free \(k\)-module \(V\), let \(E_N(V)\) be the differential
subalgebra of the absolute PD de Rham algebra
\(\Gamma(V)\otimes\Lambda(V)\) generated by the scalar unit and
all degree-zero functions of weight at least \(N\).
For \(q>0\), set
\[
 K_N^{w,q}=E_N^{w,q}(\Phi)\cap
       \bigl(dT\wedge F^{w-1,q-1}\bigr)
       \subset\Gamma(\Phi)\otimes\Lambda(\Phi).
 \tag{D-28a}
\]
There is a natural exact sequence of simplicial modules
\[
 0\longrightarrow K_N^{w,q}\longrightarrow E_N^{w,q}(\Phi)
   \longrightarrow D_N^{w,q}\longrightarrow0,
 \tag{D-28b}
\]
whose middle term is contractible. Consequently
\[
 \pi_0(D_N^{w,q})=0,\qquad
 \pi_j(D_N^{w,q})\cong\pi_{j-1}(K_N^{w,q})\quad(j\ge1).
 \tag{D-28c}
\]
In particular, the remaining eventual horizontal problem
for this construction is an eventual acyclicity problem
for the explicit intersections in (D-28a).

Higher horizontal homology does occur at finite weights:
\[
                 \pi_3(D_2^{4,1}(\Z))\cong\Z/2.
 \tag{D-28d}
\]
This finite-weight class does not refute eventual acyclicity.
\end{proposition}

\begin{proof}
In positive form degree, \(E_N(V)\) consists of sums
\(a_0\,da_1\cdots da_q\), where each \(a_i\), \(i>0\),
has weight at least \(N\), and \(a_0\) is either a scalar
or has weight at least \(N\). This differential subalgebra
is closed under the standard divided powers. Indeed every
decomposable positive even form has all higher divided
powers zero, and the divided powers of a sum expand into
products of distinct such summands. This follows by
naturality from the free exterior algebra on the displayed
one-forms. Degree-zero divided powers stay among the
generators of weight at least \(N\).

For fixed \(w>0,q\), the construction \(V\mapsto E_N^{w,q}(V)\)
is a reduced functor on free modules. Applying it to the
maps of the simplicial linear contraction of \(\Phi\)
preserves all simplicial homotopy identities: these are
identities of compositions. Thus \(E_N^{w,q}(\Phi)\) is
contractible. The terms with \(q>0,w=0\) are zero.

Projection from absolute to relative forms kills \(dT\).
Its image on \(E_N\) is all of \(D_N^q\), for \(q>0\).
The only additional coefficients allowed in \(D_N\) are
low-weight elements \(c(T)\in S\), and these can be
absorbed into the first differentiated function:
\[
 c(T)\,df_1\cdots df_q
       =d(c(T)f_1)\,df_2\cdots df_q\pmod{dT}.
 \tag{D-28e}
\]
The function \(c(T)f_1\) still has weight at least \(N\).
The kernel of absolute-to-relative projection is exactly
\(dT\wedge F^{w-1,q-1}\). This description is independent
of a splitting: two lifts of a relative exterior form
differ by a multiple of \(dT\), which its wedge with
\(dT\) kills. Degreewise splitting of the marked line
proves injectivity of this exterior inclusion.
This proves (D-28b), and its long exact homotopy sequence
proves (D-28c).

We give an exact integral matrix certificate for (D-28d).
Use cumulative coordinates \(T,X_1,\ldots,X_n\), and let
\(\alpha\) denote the multidegree counting the differential
as one occurrence of its variable. In a fixed multidegree,
the lattice of the actual image is spanned by \(a\,df\),
where \(a+f=\alpha\), the weight of \(f\) is at least \(N\),
and \(a\) is either supported on \(T\) or has weight at
least \(N\). Its coefficient at the basis one-form with
differential \(dX_j\) is
\[
 \begin{cases}
 \displaystyle\prod_{i=0}^n
           \binom{\alpha_i-\delta_{ij}}{a_i},&f_j>0,\\
 0,&f_j=0.
 \end{cases}
 \tag{D-28f}
\]
A missing variable \(X_j\) is precisely a degeneracy.
The image lattice in a multidegree missing \(X_j\) is
the identical image calculation in one fewer variable.
Therefore the normalized quotient of the image embeds
in that of the ambient forms: retain exactly the
multidegrees with \(\alpha_j\ge1\) for every \(j>0\).
This proves that the following calculation concerns
the actual integral lattice, and not the image after
reduction modulo a prime.

For \(N=2,w=4\), use lexicographic multidegrees and the
column Hermite bases from (D-28f). The normalized ranks
in simplicial dimensions \(1,2,3,4\) are \(4,12,12,4\);
all higher terms are zero. The alternating faces give
the following two matrices, with columns indexing the
source:
\begingroup\small\setlength{\arraycolsep}{2pt}
\[
d_3=\begin{pmatrix}
3&1&-1&3&-1&1&0&0&0&0&0&0\\
0&1&2&0&2&1&0&0&0&0&0&0\\
-2&-2&0&0&0&0&2&0&0&0&0&0\\
0&0&-2&0&0&0&0&2&2&0&0&0\\
0&0&0&-1&-1&1&-1&-1&1&0&0&0\\
0&0&0&0&0&-3&0&0&-3&0&0&0\\
0&1&0&0&0&0&0&0&0&2&0&0\\
0&0&1&0&0&0&0&0&0&0&1&1\\
0&0&0&0&1&0&0&0&0&-1&-1&0\\
0&0&0&0&0&1&0&0&0&0&0&-2\\
0&0&0&0&0&0&0&1&0&0&1&-1\\
0&0&0&0&0&0&0&0&1&0&0&2
\end{pmatrix},\qquad
d_4=\begin{pmatrix}
-2&0&-1&-1\\
0&-2&0&0\\
0&0&-1&-1\\
2&1&1&0\\
0&1&1&0\\
0&0&0&2\\
-2&-2&-1&-1\\
0&0&-1&1\\
0&0&0&-2\\
0&1&0&0\\
0&0&1&0\\
0&0&0&1
\end{pmatrix}.
\]
\endgroup
The matrices satisfy \(d_3d_4=0\). The last three rows
of \(d_4\) eliminate columns two through four from its
first nine rows by integral row operations. What remains
in its first column has gcd two. Its nonzero Smith
factors are consequently \(1,1,1,2\). The minor of
\(d_3\) with one-based row indices
\((1,2,3,4,5,7,8,9)\) and column indices
\((1,2,3,4,5,6,8,9)\) has determinant \(24\).
Thus \(\operatorname{rank}d_3\ge8\); the relation
\(d_3d_4=0\) and \(\operatorname{rank}d_4=4\) give the
opposite inequality.

It follows that \(H_3\) has free rank \(12-8-4=0\).
In the exact sequence
\[
0\longrightarrow H_3\longrightarrow\operatorname{coker}d_4
  \longrightarrow\operatorname{im}d_3\longrightarrow0,
\]
the last term is free, so the torsion of \(H_3\) is
exactly that of \(\operatorname{coker}d_4\), namely
\(\Z/2\). This proves (D-28d).
\end{proof}

\paragraph{Certificates and scope.}
The script \texttt{code/D\_tail\_H3\_certificate.py}
saves every normalized lattice basis, the complete
face matrices, the displayed minors, and an additional
check of all \(495\) maximal minors of \(d_4\) in
\texttt{logs/D\_tail\_H3\_certificate.json}.
The broader integral computation
\texttt{code/D\_tail\_integral\_horizontal.py}
checks degrees one through three in eleven cases
\(N=2,3\), \(N\le w\le7\). It also finds
\(\pi_3(D_3^{6,1}(\Z))=\Z/3\).
The separate script \texttt{code/D\_tail\_integral\_modp.py}
reduces the integral Hermite lattices, not merely the
defining generators, modulo \(p\). In its thirty-two
complete normalized complexes for
\(N=2,w=4,\ldots,8\) and \(N=3,w=6,7,8\), with
\(p=2,3,5,7\), the only nonzero groups occur in degrees
three and four at \((N,w,p)=(2,4,2),(3,6,3)\).
These finite calculations do not establish an eventual
bound on (D-28a). Nor does the intersection reduction
construct maps to arbitrary targets: the degree-zero
standard-tail obstruction in Statement~\ref{sa08:C13}
still applies.
